
\documentclass{article} % For LaTeX2e
\usepackage{iclr2027_conference,times}

\usepackage{hyperref}
\usepackage{url}
\usepackage{amsmath}
\usepackage{amssymb}
\usepackage{bm}
\usepackage{booktabs}
\usepackage{multirow}
\usepackage{tabularx}
\usepackage{graphicx}
\usepackage{xcolor}
\usepackage{hyperref}
\usepackage{url}
\usepackage[section]{placeins}
\usepackage{float}
\usepackage{array}
\usepackage{colortbl}
\usepackage{algorithm}
\usepackage{algpseudocode}
\usepackage{cuted}
\usepackage{caption}
\usepackage{placeins}
\usepackage{wrapfig}
\usepackage{needspace}
\setlength{\textfloatsep}{.5pc}
\setlength{\intextsep}{.5pc}
\setlength{\floatsep}{.5pc}
\usepackage{tabularx}


% ============================================================
% COLOR-BLIND-SAFE ACADEMIC PALETTE
% ============================================================

\definecolor{ICLRBlue}{RGB}{0,114,178}
\definecolor{ICLROrange}{RGB}{230,159,0}
\definecolor{ICLRGreen}{RGB}{0,158,115}
\definecolor{ICLRVermillion}{RGB}{213,94,0}
\definecolor{ICLRPurple}{RGB}{204,121,167}
\definecolor{ICLRSky}{RGB}{86,180,233}
\definecolor{ICLRGray}{RGB}{90,90,90}
\definecolor{ICLRLightGray}{RGB}{225,225,225}
\definecolor{ICLRVeryLightGray}{RGB}{246,246,246}
\definecolor{RQBlue}{RGB}{78,121,167}
\definecolor{RQTeal}{RGB}{102,194,165}
\definecolor{RQOrange}{RGB}{242,142,44}
\definecolor{RQPurple}{RGB}{141,160,203}
\definecolor{RQGray}{RGB}{140,140,140}
\definecolor{RQGrid}{RGB}{225,225,225}
\definecolor{purpleA}{RGB}{235,235,247}
\definecolor{purpleB}{RGB}{214,213,237}
\definecolor{purpleC}{RGB}{184,183,220}
\definecolor{purpleD}{RGB}{151,145,201}
\definecolor{purpleE}{RGB}{117,105,174}

\definecolor{purpleEdgeA}{RGB}{180,180,213}
\definecolor{purpleEdgeB}{RGB}{161,159,203}
\definecolor{purpleEdgeC}{RGB}{132,128,185}
\definecolor{purpleEdgeD}{RGB}{102,92,157}
\definecolor{purpleEdgeE}{RGB}{79,67,132}



\title{
Confidence-Gated Mechanistic\\
Bridge Retrieval for Biomedical\\
Knowledge Graph Reasoning
}

% Authors must not appear in the submitted version. They should be hidden
% as long as the \iclrfinalcopy macro remains commented out below.
% Non-anonymous submissions will be rejected without review.

\author{
Shiyu Xi
\And
Yushun Dong
}

% The \author macro works with any number of authors. There are two commands
% used to separate the names and addresses of multiple authors: \And and \AND.
%
% Using \And between authors leaves it to \LaTeX{} to determine where to break
% the lines. Using \AND forces a linebreak at that point. So, if \LaTeX{}
% puts 3 of 4 authors names on the first line, and the last on the second
% line, try using \AND instead of \And before the third author name.

% Blue highlights revised data/results/conclusions; general prose is unmarked.
\newcommand{\rev}[1]{{\color{blue}#1}}
\newcommand{\fix}{\marginpar{FIX}}
\newcommand{\new}{\marginpar{NEW}}

%\iclrfinalcopy % Uncomment for camera-ready version, but NOT for submission.
\begin{document} \raggedbottom


\maketitle

\maketitle
% ============================================================
% ABSTRACT
% ============================================================
\begin{abstract}
Biomedical knowledge-graph reasoning requires recovering intermediate mechanistic evidence rather than only connecting a source and target. Semantic guidance can improve graph retrieval, but its value is conditional: useful candidates may be missing, semantically plausible candidates may provide limited downstream retrieval value, and the guidance strength may vary across queries. We introduce \textbf{Confidence-Gated Mechanistic Bridge Retrieval}, a graph-constrained framework that separates these challenges into candidate coverage, retrieval utility, and query-specific influence. The method constructs graph-linked Bridges from complementary biological views, estimates retrieval-aware reliability from downstream supervision, and adaptively controls their influence during graph-valid search. On DrugMechDB \rev{with supplied reference-aligned candidates, Full improves Path Success from 27.02\% to 39.52\% over graph-only retrieval; a separate automatic-proposal evaluation reaches 35.08\%.} Controlled analyses distinguish candidate coverage, downstream Bridge utility, and query-specific influence; \rev{the incremental benefits of learned selection and control are smaller and depend on the comparison.} These results support a failure-oriented view of semantic graph guidance, in which availability, downstream value, and query-specific influence should be modeled separately. An anonymous repository for code and artifacts is available at
\url{https://anonymous.4open.science/r/iclr2027-bridge/}.
\end{abstract}
% ============================================================
% 1 INTRODUCTION
% ============================================================
\section{Introduction}
\label{sec:introduction}

Retrieval-augmented generation (RAG) gives language models access to external
evidence during reasoning. When the external knowledge is represented as a
graph, retrieval can also exploit explicit entities, relations, and multi-hop
connectivity. Recent graph-based RAG methods use graph traversal,
query-conditioned subgraphs, hierarchical structures, and learned retrievers
to improve structured evidence access
\citep{sun2024tog,ma2025tog2,he2024gretriever,
mavromatis2025gnnrag,li2025subgraphrag,zhang2025sirerag,
zhuang2026linearrag,luo2026greasoner}.
This setting is especially relevant to biomedical reasoning, where knowledge
spans drugs, genes, proteins, pathways, biological processes, diseases, and
phenotypes
\citep{himmelstein2017hetionet,chandak2023primekg}.
Mechanistic questions place an additional demand on retrieval. They require
not only relevant facts, but also the intermediate biological relations that
explain how a source entity leads to a target
\citep{gonzalez2023drugmechdb,callaghan2023biothings,joy2026bterag}.

Existing methods improve graph retrieval by selecting relations or entities,
ranking query-relevant graph regions, and adapting graph organization or
search to the query
\citep{zhou2026qafd,zhao2026toporag}.
Recent work studies graph constraints, relation-aware retrieval, and medical evidence
\citep{luo2025gcr,wu2025medgraphrag,um2026rporag}.
These advances improve which evidence is retrieved. A different issue arises
when intermediate semantic information is used to guide graph search.
Semantic guidance can direct retrieval toward meaningful mechanisms, but it
can also mislead search when the proposed mechanism is incomplete or
irrelevant. Its usefulness may also differ across queries. Adaptive RAG work
has shown that retrieval behavior can depend on confidence and query
properties
\citep{jiang2023flare,asai2024selfrag,jeong2024adaptiverag}, but this issue
remains less understood when semantic guidance operates inside
graph-constrained retrieval.

This setting raises a specific design question: how should an intermediate
mechanistic hypothesis enter graph retrieval? Its contribution depends on
three decisions made at different points in the pipeline. A relevant
intermediate entity must first be available among the candidates; among those
candidates, semantic plausibility need not imply value for retrieving the
reference mechanism; and even useful guidance need not have the same influence
for every query. These are decisions about coverage, downstream utility, and
influence, not different names for a single relevance score.

We therefore introduce Confidence-Gated Mechanistic Bridge Retrieval, a graph-constrained framework built around these three decisions. A Mechanistic Bridge is a graph-linked intermediate hypothesis that provides a semantic signal during path search. The framework constructs candidate Bridges from complementary biological views to broaden access to intermediate mechanisms. It then estimates each candidate's reliability from its marginal contribution to downstream retrieval, rather than from textual similarity alone. Finally, a query-level controller determines how strongly the selected Bridges affect path ranking. All search steps remain graph-valid: Bridges prioritize existing paths but do not add edges to the graph. This decomposition also gives each component a clear empirical role: candidate coverage can be measured before search, Bridge reliability can be evaluated under matched candidate and search conditions, and query-specific influence can be compared against constant or reassigned guidance strengths. The resulting experiments therefore \rev{test how each retrieval decision affects mechanistic evidence recovery}, without requiring every stage to produce a large, monotonic gain in aggregate endpoint accuracy. In particular, matched retrieval conditions separate Bridge utility from candidate access, while reassigned guidance strengths distinguish query-specific influence from average semantic guidance.


Our contributions are:
\begin{itemize}
    \item \textbf{Separating the decisions behind semantic guidance.}
    We identify a three-part decomposition of graph guidance:
    whether intermediate candidates are available, whether they improve
    downstream retrieval, and how strongly they should influence each
    query. This makes the three sources of failure independently testable.

    \item \textbf{Retrieval-aware Mechanistic Bridge guidance.}
    We introduce a graph-constrained framework that constructs
    Bridges from complementary biological views, estimates their
    retrieval-derived reliability, and controls their influence per
    query. Bridges alter path ranking while all retrieved paths remain
    valid in the original graph.


    \item \textbf{Isolating the effects of each retrieval decision.}
    Matched interventions separate candidate access, Bridge usefulness, and query-specific influence. They show broader automatic candidate coverage and informative usefulness scores, while quantifying \rev{the limited incremental gains of selection and query-specific control.}



\end{itemize}


\section{Preliminaries}
\label{sec:preliminaries}

% ============================================================

% 2.1 TASK SETUP
% ============================================================

\subsection{Task Setup}
\label{sec:task_setup}

We represent a biomedical knowledge graph as $G=(V,E,\mathcal{R})$, where
$V=\{v_1,\ldots,v_n\}$ denotes the set of biomedical entities,
$\mathcal{R}$ denotes the relation types, and
$E\subseteq V\times\mathcal{R}\times V$ denotes directed biomedical
relations \citep{himmelstein2017hetionet,chandak2023primekg}. Each graph fact
is represented as a triple $(v_i,r,v_j)\in E$, and explicit graph structure
supports retrieval and reasoning over entities, relations, paths, and
query-conditioned graph regions
\citep{das2018minerva,sun2024tog,ma2025tog2,yuan2025kgrag,
edge2024graphrag,he2024gretriever}. Recent work further explores
knowledge-graph-guided retrieval, differentiable KG retrieval, multi-strategy
graph retrieval, structured-data reasoning, and the broader integration of
large language models with knowledge graphs
\citep{zhu2025kg2rag,gao2025drag,mavromatis2025byokg,
jiang2023structgpt,ma2025llmkgqa}. Together, these settings motivate retrieval
procedures that exploit explicit graph structure while preserving the
distinction between graph evidence and auxiliary semantic guidance.

Given a biomedical mechanistic question $q$, we identify a source entity
$e_s\in V$, a target entity $e_t\in V$, and a reasoning intent $I_q$, and
denote the resulting structured query as $Q=(q,e_s,e_t,I_q)$. Our primary
setting is endpoint-conditioned retrieval, in which both endpoints are
available at retrieval time and the target therefore defines the retrieval
endpoint rather than an entity to be predicted. A graph-valid path is written
as $p=(e_s,r_1,v_1,\ldots,r_L,e_t)$, where $L$ denotes the number of
relations, and $\mathcal{P}(e_s,e_t)$ denotes the set of graph-valid candidate
paths between the endpoints
\citep{das2018minerva,lin2018multihop,sun2019pullnet,
gutierrez2025hipporag2,guo2024lightrag}. The retrieved evidence
$G_E\subseteq G$ may contain one or more supporting paths together with their
associated entities and relations. Throughout the paper, semantic information
may alter the ranking of graph-valid candidates, but it does not create
virtual edges or modify graph connectivity.
% ============================================================
% 2.2 PROBLEM STATEMENT
% ============================================================
\subsection{Problem Statement}
\label{sec:problem_statement}

We study retrieval of intermediate mechanistic evidence between known biomedical endpoints. Unlike relation prediction or answer recovery~\citep{das2018minerva,lin2018multihop,sun2019pullnet}, this setting evaluates the biological entities and relations along a retrieved connection. DrugMechDB provides such ordered intermediate relations as curated reference mechanisms~\citep{gonzalez2023drugmechdb}. They trace the biological relations linking the source and target entities.

\noindent\textbf{Definition 1 (Mechanistic Bridge).} For a query $Q=(q,e_s,e_t,I_q)$ on $G$, a Mechanistic Bridge $h$ is a query-relevant hypothesis about an intermediate biomedical entity that can be linked to a node of $G$. A Bridge identifies a candidate waypoint for retrieval; it is neither a graph edge nor evidence that a complete mechanism exists. Let $H_0(Q)$ be the set of grounded Bridge candidates available for $Q$.

\noindent\textbf{Problem 1 (Mechanistic Evidence Retrieval).} Given $G$, $Q$, and a retrieval budget $B$, return an evidence subgraph $G_E\subseteq G$ comprising graph-valid paths from $e_s$ to $e_t$ within $B$. The objective is to recover intermediate entities and relations relevant to the mechanism in $q$, rather than endpoint connectivity alone. When reference mechanisms are available, we assess this objective using path success and mechanistic edge recall. Bridge candidates may guide which valid paths are explored and ranked, but cannot change the feasible graph paths.

This formulation separates three choices in semantic guidance: which intermediate hypotheses are available, which of them improve downstream retrieval, and how much their guidance should affect a particular query. The method addresses these choices through candidate construction, retrieval-aware selection, and query-specific influence control under a shared retrieval budget.
% ============================================================
% 3 AN EXPLORATORY STUDY
% ============================================================
\section{An Exploratory Study}
\label{sec:exploratory}



Graph grounding rules out invented edges, but it does not establish that a grounded intermediate hypothesis will lead search toward the desired mechanism. A Bridge can name a real biological entity and still favor a plausible path that misses the reference relations, despite remaining fully valid in the graph. This creates a question left open by the task formulation: when graph connectivity and the search budget are held fixed, can changing only the intermediate guidance turn a recoverable mechanism into a missed one? A second question follows even if the Bridge is informative: should it exert the same influence on every query?

\begin{figure}[H]
    \centering
    \setlength{\abovecaptionskip}{.5pc}
    \setlength{\belowcaptionskip}{0pt}

    \includegraphics[
        width=0.82\linewidth,
        trim=7 180 24 68,
        clip
    ]{figure 1.pdf}

    \caption{
    Semantic guidance is conditional: degrading Bridge quality reduces retrieval benefit, while fixed guidance produces heterogeneous responses across queries.
    }
    \label{fig:exploratory_bridge}

\end{figure}

We examine these questions through a controlled corruption study. Starting from reference-aligned Clean Bridges, we progressively replace intermediate entities with type-compatible distractors drawn from the same biomedical graph vocabulary, preserving the form of the guidance while reducing its mechanistic accuracy. Random contains only distractor guidance, and No Bridge provides the graph-only reference. Across these conditions, the graph, source--target pair, and search budget are matched. The comparison therefore attributes changes in retrieved evidence to the supplied Bridge information rather than to altered connectivity or search capacity.

\begin{figure}[!t]
    \centering
    \setlength{\abovecaptionskip}{0pt}
    \setlength{\belowcaptionskip}{0pt}

    \includegraphics[
        width=\linewidth,
        trim=0 30 0 8,
        clip
    ]{figure2.pdf}

    \caption{Overview of the three-stage retrieval framework.
    Multi-view Bridge construction forms $H_0$, retrieval-aware selection
    yields $H^*$, and confidence-gated graph retrieval uses $g$ to control
    query-specific semantic influence while preserving graph-valid search.
    }
    \label{fig:framework_overview}

\end{figure}

Figure~\ref{fig:exploratory_bridge} reveals two distinct behaviors. First,
guidance utility deteriorates with Bridge quality: Clean Bridges benefit more
retrieval instances, whereas progressive corruption shifts queries toward
neutral or harmful outcomes. Because the graph and search budget are matched,
this change isolates the effect of the supplied intermediate information and
shows that unreliable Bridges can redirect search away from otherwise
recoverable evidence. Second, guidance effects remain heterogeneous across
queries even under the same fixed coefficient: some instances benefit from the
supplied semantic information, whereas others receive little or negative
benefit. This intervention varies supplied guidance rather than candidate
generation, so it does not measure whether useful Bridges are available in
realistic proposals. Instead, it shows why grounded candidates must be assessed
for downstream utility and why their influence should vary across queries.
% ============================================================
% 4 METHODOLOGY
% ============================================================

\section{Confidence-Gated Mechanistic Bridge Retrieval}
\label{sec:method}
% ============================================================
% 4.1 OVERVIEW
% ============================================================
\subsection{Overview}
\label{sec:method_overview}

We present the overall architecture in Figure~\ref{fig:framework_overview}. Given $Q$ and $G$, the framework implements the three guidance choices introduced in (Section~\ref{sec:problem_statement}). Stage I constructs a graph-linked candidate pool $H_0$ from complementary molecular, cellular, and system-level views to address candidate coverage (Section~\ref{sec:bridge_construction}); Stage II estimates retrieval-aware Bridge reliability and reduces $H_0$ to a reliable and complementary subset $H^*$ to address downstream utility (Section~\ref{sec:bridge_reliability}); and Stage III maps query-level reliability to a continuous guidance strength $g$ that controls the influence of $H^*$ during graph-constrained path ranking (Section~\ref{sec:confidence_gate}). The resulting search proceeds from candidate availability, to candidate usefulness, to query-specific influence, yielding an evidence subgraph $G_E$ composed only of graph-valid supporting paths. Semantic information changes only the ranking of valid graph candidates and never alters graph connectivity or introduces non-graph edges.
% ============================================================
% 4.2 MECHANISTIC BRIDGE CONSTRUCTION
% ============================================================
\subsection{Mechanistic Bridge Construction}
\label{sec:bridge_construction}
\textbf{Multi-view candidate construction.} Stage I addresses candidate-coverage failure by constructing graph-linked Mechanistic Bridges from complementary biological views. We use three view groups, $\mathcal{V}=\{\text{molecular},\text{cellular},\text{system}\}$, so that candidate generation is not restricted to a single abstraction level. In our implementation, the molecular view covers entity types such as proteins, chemical substances, molecular activities, and gene families. The cellular view covers cells, cellular components, macromolecular complexes, and gross anatomical structures. The system-level view covers biological processes, pathways, phenotypic features, and organism-level entities. For each view $v\in\mathcal{V}$, a view-specific proposer produces a candidate set $\tilde{H}^{(v)}=\mathrm{Prop}_v(Q)$. We combine the candidates as $\tilde{H}=\bigcup_{v\in\mathcal{V}}\tilde{H}^{(v)}$, then normalize surface forms and remove duplicates before graph grounding. The candidate pool is $H_0=\mathrm{Link}(\mathrm{Dedup}(\mathrm{Normalize}(\tilde{H})),G)$. Only candidates that can be mapped to entities already present in $G$ are retained. Stage I therefore expands the set of available intermediate hypotheses without adding new graph edges or changing graph topology.

For automatic evaluation, $\mathrm{Prop}_v$ enumerates two forward rounds from $e_s$, followed by two reverse rounds from $e_t$, with a shared visited set, endpoint exclusion and the above type filters. Canonical graph IDs provide grounding and deduplication. A fixed lexical/adjacency/anti-hub/relation-quality rule ranks proposals. Inputs are known endpoints, task intent and the fixed graph; neither the intermediate answer nor query-specific reference membership is read by the proposer. Appendix~\ref{app:implementation} specifies the rule and distinguishes automatic pools from supplied Clean fixtures.

\noindent\textbf{Controlled view partition.} 
The reference-aligned controlled partition uses the same biological-view grouping to isolate candidate access; Section~\ref{sec:stage1_eval} separately evaluates the executable automatic proposer. In this setting, held-out Clean candidates are assigned to molecular, cellular, or system-level groups using the same entity-type mapping. Retrieval-aware selection and adaptive influence control are disabled, and all retained Bridges are assigned equal confidence under a fixed semantic guidance strength. The multi-view condition uses the union of the three grounded view-specific candidate sets under the same retrieval procedure. This controlled construction separates the effect of candidate availability from later reliability and influence mechanisms. It provides the methodological basis for evaluating whether Stage I broadens access to useful intermediate evidence.
% ============================================================
% 4.3 RETRIEVAL-AWARE BRIDGE RELIABILITY
% ============================================================
\subsection{Retrieval-Aware Bridge Reliability}
\label{sec:bridge_reliability}

\noindent\textbf{Bridge representation.} Stage II addresses utility mismatch by estimating whether each candidate Bridge can improve downstream mechanistic retrieval, rather than whether it is only semantically related to the query. For each $h\in H_0$, we construct a representation $x_h=[x_h^{\mathrm{sem}};x_h^{\mathrm{graph}}]$. The semantic component is built from the task label, query text, Bridge name, and Bridge type using word-level TF--IDF features with 1--2 grams and character-level TF--IDF features with 3--5 grams. The graph component describes pre-search structural context. It includes log-transformed in-degree, out-degree, and total degree, source adjacency, source out-degree, target in-degree, normalized incidence over the 30 most frequent relation types, and node-type indicators. Text vocabularies and numeric standardization statistics are fitted on the training split only. We do not use reference-path membership, target-path oracle information, or post-retrieval outcomes as input features.

\noindent\textbf{Retrieval-derived reliability supervision.} We define Bridge usefulness by its marginal effect on retrieval utility $M(H)=\omega\,\mathrm{PathSuccess}(H)+(1-\omega)\,\mathrm{EdgeRecall}(H)$. The default weight is $\omega=0.7$; its label sensitivity is examined in Appendix~\ref{app:search_parameter_sensitivity}. The marginal contribution of candidate $h$ is $\Delta_h=M(H\cup\{h\})-M(H)$, and a candidate is labeled useful when $\Delta_h>0$. This supervision therefore reflects whether adding $h$ improves retrieval under matched graph and search conditions. It does not use semantic plausibility as the target. Given $x_h$, we estimate a usefulness score as $c_h=\sigma(w_c^\top x_h+b_c)$. In implementation, the estimator is a class-balanced logistic model trained only on the training split. The resulting score $c_h\in[0,1]$ is used as retrieval-aware Bridge reliability. Class balancing means this score is \rev{not assumed to be a calibrated population probability, nor does positive-benefit classification guarantee ranking by gain magnitude.}

\noindent\textbf{Reliability-aware Bridge selection.} High-confidence candidates can still provide redundant guidance, so Stage II does not select Bridges using $c_h$ alone. Given the currently selected set $H$, each remaining candidate is scored by $\mathrm{Gain}(h\mid H)=c_h-\mu\,\mathrm{Red}(h,H)+\beta\,\mathrm{TypeBonus}(h,H)$. The redundancy term penalizes overlap with already selected Bridges, while the type bonus favors complementary mechanistic types. Starting from an empty set, candidates are added greedily according to this gain until the Bridge budget is reached. This produces the selected set $H^*\subseteq H_0$, which is passed to Stage III for query-specific influence control. \rev{The archived main-table VAL/TEST pools contain at most $K=3$ nodes, so they do not test subset filtering.} We therefore distinguish confidence weighting from the complete greedy rule, and evaluate the latter on oversized automatic pools under the stated gain rule (Appendix~\ref{app:implementation}).
% ============================================================
% 4.4 CONFIDENCE-GATED GRAPH RETRIEVAL
% ============================================================
\subsection{Confidence-Gated Graph Retrieval}
\label{sec:confidence_gate}

\paragraph{Query-specific influence control.}
Stage III addresses influence mismatch by determining how strongly the selected semantic guidance should affect each query during retrieval. Given the selected Bridge set $H^*$, we construct a query-level representation $x_Q$ from the task label, query text, Bridge names and types, together with graph information available before search. Structural features include normalized candidate-set size, log source out-degree, log target in-degree, candidate-type proportions, source-adjacency rate, and mean log candidate degree. These features describe the query and its semantic context without using shortest-path or post-retrieval information. Query-level supervision is derived directly from retrieval outcomes. For each training instance, we define $z=\mathbb{I}\!\left[\max_{g'>0}M(g';q,H^*)>M(0;q,H^*)\right]$, where $M(g';q,H^*)$ is the retrieval utility under positive semantic guidance and $M(0;q,H^*)$ is the corresponding graph-only utility. Thus, $z=1$ indicates that positive guidance can improve retrieval for the query. We estimate query-level reliability as $p_{\mathrm{rel}}=\sigma(w_g^\top x_Q+b_g)$ and map it to a continuous strength $g=g_{\max}p_{\mathrm{rel}}^\gamma$. The parameters $g_{\max}$ and $\gamma$ are selected using validation retrieval utility. Lower reliability therefore reduces semantic influence without imposing a binary guidance decision.
\paragraph{Graph--Bridge path scoring.}
For each graph-valid path $p$, we combine structural transitions with agreement between $p$ and the selected Bridges. The graph score $S_{\mathrm{graph}}(p)$ averages local scores $s_{\mathrm{local}}(v_i,r_i)=\lambda q_{\mathrm{rel}}(r_i)+(1-\lambda)(1-\tilde d(v_i))$. Relation quality favors mechanistic edges, while the degree term discounts hubs. The Bridge score is $S_{\mathrm{bridge}}(p,H^*)=\operatorname{clip}_{[0,1]}[\rho(p)C(p,H^*)]$. The summary $C$ combines Bridge coverage, mean reliability, current endpoint reliability, and recent Bridge hits. Coverage has the largest contribution; the endpoint and recency terms receive equal, smaller contributions. The factor $\rho(p)$ rises as the path reaches more distinct selected Bridges. It gives repeated mechanistic support more weight than one incidental hit. This distinction matters when several paths contain a Bridge but only some reach multiple selected mechanisms. Appendix~\ref{app:implementation} specifies the fixed scorer values, and Appendix~\ref{app:search_parameter_sensitivity} reports their local sensitivity.

\paragraph{Graph-constrained retrieval.}
The path-ranking score is $S(p)=(1-g)S_{\mathrm{graph}}(p)+gS_{\mathrm{bridge}}(p,H^*)$. The gate changes how graph and Bridge evidence rank valid paths. Search begins at the source and expands only edges present in $G$. A fixed beam width and expansion budget limit this process. After each step, we retain the highest-scoring frontier states. We exclude repeated nodes to prevent cycles. Search stops at the target or when the budget is exhausted. The retained path then defines the evidence subgraph $G_E$. The graph and feasible paths remain fixed as the gate changes. The controller does not differentiate through beam search and receives no answer-generation loss. Bridge guidance can prioritize one valid path over another, but it cannot add edges or make an unreachable target reachable. Matched controls use identical graphs and search budgets. They also retain the same reachable candidates and the same stopping rule across guidance settings. This separation matters for interpreting the ablations: improved retrieval cannot be attributed to expanded reachability or a more permissive search. The resulting differences in recovered mechanistic evidence can therefore be attributed to path ranking under a common graph and search budget.

% 5 EXPERIMENTAL EVALUATION
% ============================================================
\section{Experimental Evaluation}
\label{sec:experiments}

We evaluate the framework through four research questions.
\textbf{RQ1:} How do Bridge guidance, confidence weighting, and
query-specific influence affect retrieval with fixed candidate sets?
\textbf{RQ2:} Does multi-view Bridge construction reduce candidate-coverage
failure?
\textbf{RQ3:} Can retrieval-aware supervision distinguish downstream-useful
Bridges from merely plausible candidates?
\textbf{RQ4:} Does query-specific influence control outperform uniform
semantic guidance by matching guidance strength to each query?
Unless otherwise stated, model selection and hyperparameter tuning use training
and validation data only, with held-out test sets reserved for final evaluation.

\subsection{Experimental Setup}
\label{sec:setup}

\textbf{Datasets.} We evaluate the core framework on DrugMechDB
\citep{gonzalez2023drugmechdb}, using the BTE-RAG benchmark family
\citep{joy2026bterag} of DMDB-Gene, DMDB-Metabolite, and DMDB-Process.
Gene and Metabolite evaluate recovery of intermediate mechanistic entities
from drug--disease contexts, while Process is reformulated as
drug+disease $\rightarrow$ biological process to avoid answer leakage.
EpiKG, released with EpiGraph \citep{dai2026epigraph}, provides an
independent biomedical graph used to test whether retrieval-aware Bridge
reliability remains informative under a different graph schema and candidate
distribution. This setting complements DrugMechDB by changing the graph
structure and supervision while preserving the evaluation of graph-grounded
intermediate guidance. Additional transfer evaluations on OpenBioLink
\citep{breit2020openbiolink}, BioHopR \citep{kim2025biohopr}, and
ComplexWebQuestions (CWQ) \citep{talmor2018web} are reported in
Appendix~D. These experiments probe whether the proposed guidance principles
remain informative across alternative biomedical graphs, multi-hop retrieval
settings, and a general-domain graph. They assess transfer beyond the primary mechanistic retrieval benchmark.

\noindent\textbf{DrugMechDB protocol.} The manifest has \rev{1,354/239/248} TRAIN/VAL/TEST task queries from \rev{1,120/202/218} mechanism records. No record ID or its candidate/corruption variants crosses splits. Different records can share endpoint pairs: \rev{9 test queries} have a pair seen in TRAIN or VAL, so this is not strict unseen-pair evaluation. The fixed aggregate graph includes held-out mechanism edges; evaluation is transductive. Main Clean fixtures are reference-aligned (\rev{692 test candidate records}, with the answer present in every query pool); automatic proposal is reported separately.

\paragraph{\textbf{Evaluation Metrics.}}
We use benchmark-appropriate metrics while preserving the same evaluation
principle within each task. DrugMechDB reports Path Success (PS), Edge Recall
(ER), and Answer Hit (AH): PS measures recovery of a graph-valid
source--target path, ER measures recovery of directed edges in the reference
mechanism, and AH indicates whether the benchmark answer entity appears in the
retrieved evidence. Graph-recovery evaluations use the corresponding
path- or edge-recovery metrics, while target-agnostic ranked-answer settings
use Hits@K and MRR under their native evaluation protocols. Bridge-level and
query-level reliability estimators are evaluated using AUROC, AUPRC, and Brier
score. DrugMechDB paired comparisons use 5,000 bootstrap draws of mechanism-record IDs, retaining all tasks and candidate variants of a record together, with pointwise 95\% intervals; shuffle tests additionally report multiplicity adjustment. Scores are compared only within the same benchmark because task formulations and metric scales differ.

\textbf{Baselines.} Graph Only uses structural scores. In the supplied Clean main-table pools, Semantic+Fixed assigns every Bridge confidence 1 with $g=0.3$, while Utility+Fixed uses learned confidences at the same $g$. \rev{These rows retain the same nodes and test confidence weighting, not subset selection.} Full adds the original continuous gate. Stage-III controls fix these same nodes and confidences, then compare validation-selected global, validation-mean and shuffled gates. A compact learned-path scorer additionally matches supervision source and search budget (Appendix~\ref{app:implementation}); its path-level labels differ from the gate's query-level labels. Oracle variants are non-deployable.
% ============================================================
% 5.2 OVERALL END-TO-END RETRIEVAL CONTEXT
% ============================================================

\subsection{Overall End-to-End Retrieval Context}
\label{sec:overall}

\noindent\textbf{Full-pipeline context.}
To answer RQ1, Table~\ref{tab:overall_gain} compares retrieval on supplied reference-aligned Clean candidates under the same graph and search budgets. It provides conditional downstream context, rather than evidence for automatic proposal or complete subset filtering. Section~\ref{sec:stage1_eval} reports automatic proposal, Section~\ref{sec:stage2_eval} isolates usefulness, and Section~\ref{sec:stage3} uses exactly this Table~\ref{tab:overall_gain} Full configuration for influence controls.

\begin{wraptable}{r}{0.46\textwidth}
\vspace{-5pt}
\centering
\caption{Retrieval on supplied reference-aligned Clean candidates (\%).}
\label{tab:overall_gain}

\vspace{-3pt}
\small
\setlength{\tabcolsep}{3.8pt}
\renewcommand{\arraystretch}{1.05}

\begin{tabular}{lrrrr}
\toprule
\textbf{Method} &
\textbf{PS} &
\textbf{ER} &
\textbf{AH} &
\textbf{$\Delta$ER} \\
\midrule
Graph Only
& 27.02
& 33.97
& 70.56
& -- \\

Semantic+Fixed
& 39.11
& 58.56
& 91.53
& +24.59 \\

Utility+Fixed
& 39.52
& 58.94
& 91.13
& +24.97 \\

Full
& \textbf{39.52}
& \textbf{59.46}
& \textbf{91.94}
& \textbf{+25.49} \\
\bottomrule
\end{tabular}

\vspace{-5pt}
\end{wraptable}
\noindent\textbf{Overall observations.} Bridge guidance accounts for most of the change: Semantic+Fixed already reaches ER 58.56\%, versus 59.46\% for Full. \rev{Full minus Semantic+Fixed gives ER $+0.90$ percentage points (paired 95\% CI $[0.23,1.69]$) and utility $+0.00552$ ($[0.00118,0.01275]$); PS $+0.40$ points has interval $[0.00,1.24]$. These small combined weighting/control gains do not establish each learned component's independent advantage.} \rev{A matched compact learned-path scorer reaches PS/ER 35.48/56.93\%; Full exceeds it by 4.03/2.53 points with intervals $[1.57,6.94]$/$[0.62,4.46]$. Fixed learned-confidence guidance also exceeds this baseline, so the difference cannot all be credited to the gate.} This comparison covers one numeric linear path scorer, not all learned-search methods. Scorer interventions in Appendix~\ref{app:search_parameter_sensitivity} retain the original graph and candidate access.
% ============================================================
% 5.3 COMPLEMENTARITY AND COVERAGE OF MULTI-VIEW MECHANISTIC BRIDGES
% ============================================================

\subsection{Complementarity and Coverage of Multi-View Mechanistic Bridges}
\label{sec:stage1_eval}

\noindent\textbf{Coverage and controlled retrieval.}
To answer RQ2, we evaluate the automatic graph-native proposer with reliability selection and adaptive influence disabled. Useful-candidate coverage means that at least one candidate, evaluated alone with confidence 1 and $g=0.3$, improves utility over Graph Only. \rev{The multi-view union covers 152/248 queries (61.29\%), versus 100/248 (40.32\%) for the strongest single view.} This measures candidate availability before final-budget truncation, not uniform fixed-$K$ superiority. Figure~\ref{fig:stage1_multiview} shows view-level complementarity using the same manifest-endpoint candidate pools as the automatic retrieval evaluation.

\noindent\textbf{Stage-I observations.}
With fixed equal-confidence guidance and at most three automatically proposed Bridges, \rev{PS/ER/AH are 35.08/38.73/73.39\%, versus Graph Only 27.02/33.97/70.56\%. Paired PS/ER gain intervals are $[4.40,12.11]$/$[2.88,6.84]$ percentage points.} Useful guidance therefore survives imperfect automatic proposal. The strongest view varies by task, and the multi-view advantage shrinks when final Bridge capacity is matched: \rev{Stage I broadens access rather than guaranteeing the best fixed-budget retrieval.} Automatic results remain distinct from supplied-Clean retrieval.

\begin{figure}[!ht]

    \centering
    \setlength{\abovecaptionskip}{.5pc}
    \setlength{\belowcaptionskip}{0pt}
    \includegraphics[
        width=0.98\linewidth,
        keepaspectratio
    ]{figure3.pdf}

    \caption{
Stage-I candidate availability across Gene, Metabolite, and Process queries, using the same normalized endpoint pools as automatic retrieval. Labels give the mean pool size and fraction with at least three candidates; the union contains \rev{14,029} query--candidate records across 248 queries.
}
    \label{fig:stage1_multiview}

\end{figure}
% ============================================================
% 5.4 RETRIEVAL UTILITY OF LEARNED BRIDGE RELIABILITY
% ============================================================

\subsection{Retrieval Utility of Learned Bridge Reliability}
\label{sec:stage2_eval}

\paragraph{Usefulness discrimination and controlled selection.}
To answer RQ3, we compare Lexical Similarity, a local Graph Heuristic, and
retrieval-aware reliability on the same held-out Bridge candidates. We evaluate
usefulness discrimination on the full and Clean candidate sets, followed by a
matched single-Bridge intervention that fixes candidate access, search conditions,
and semantic influence while disabling Stage~III. We further evaluate the same
Stage-II reliability formulation on EpiKG to test whether the learned utility signal
remains informative under a different biomedical graph schema and candidate
distribution. Because the released EpiBench questions are only partially grounded
in the released graph, this evaluation uses a frozen graph-grounded masked-edge
setting with disjoint source--target pairs across train, validation, and test splits.
Results are summarized in Tables~\ref{tab:stage2} and~\ref{tab:epikg}.

\begin{table}[H]
\centering
\vspace{-2mm}
\caption{Stage-II Bridge usefulness and controlled retrieval.}
\label{tab:stage2}
\vspace{-1mm}

\resizebox{\linewidth}{!}{
\begin{tabular}{@{}lccccccc@{}}
\toprule
&
\multicolumn{4}{c}{\textbf{Bridge usefulness}}
&
\multicolumn{3}{c}{\textbf{Controlled retrieval (\%)}} \\
\cmidrule(lr){2-5}
\cmidrule(lr){6-8}

\textbf{Method}
&
\multicolumn{2}{c}{\textbf{All conditions}}
&
\multicolumn{2}{c}{\textbf{Clean}}
&
\multicolumn{3}{c}{\textbf{Matched intervention}} \\

\cmidrule(lr){2-3}
\cmidrule(lr){4-5}
\cmidrule(lr){6-8}

&
\textbf{AUROC}
&
\textbf{AUPRC}
&
\textbf{AUROC}
&
\textbf{AUPRC}
&
\textbf{PS}
&
\textbf{ER}
&
\textbf{AH} \\
\midrule

Lexical Similarity
& 0.594 & 0.166 & 0.515 & 0.257
& \textbf{29.92} & 38.19 & 75.32 \\

Graph Heuristic
& 0.391 & 0.089 & 0.509 & 0.238
& 27.66 & 36.15 & 72.90 \\

Retrieval-aware Reliability
& \textbf{0.795} & \textbf{0.393}
& \textbf{0.815} & \textbf{0.655}
& 29.19 & \textbf{39.92} & \textbf{76.61} \\

\bottomrule
\end{tabular}
}

\vspace{-2mm}
\end{table}

\vspace{-0.25em}

\begin{table}[!t]
\centering
\caption{Stage-II reliability transfer on EpiKG.}
\label{tab:epikg}

\vspace{-0.65em}

\footnotesize
\setlength{\tabcolsep}{3.8pt}
\renewcommand{\arraystretch}{0.88}

\begin{tabular*}{\textwidth}{
@{\extracolsep{\fill}}
lc
lc
@{}
}
\toprule

\multicolumn{2}{c}{\textbf{EpiBench grounding}}
&
\multicolumn{2}{c}{\textbf{EpiKG Bridge reliability}}
\\[-0.18em]

\cmidrule(lr){1-2}
\cmidrule(lr){3-4}

\textbf{Criterion}
& \textbf{Value}
& \textbf{Metric}
& \textbf{Value}
\\[-0.12em]

\midrule

Questions audited
& 1,000
& Useful-Bridge prevalence
& 1.05\%
\\[-0.10em]

Correct option exactly mapped
& 10.9\%
& AUROC
& 0.881
\\[-0.10em]

All four options exactly mapped
& 4.3\%
& AUPRC
& 0.104
\\[-0.10em]

At least one option exactly mapped
& 32.8\%
& Random AUPRC reference
& 0.0105
\\[-0.10em]

Mean mapped options per question
& 0.705
& AUPRC / prevalence
& 9.9$\times$
\\[-0.05em]

\bottomrule
\end{tabular*}

\vspace{-0.65em}
\end{table}

\vspace{-0.35em}


\paragraph{Stage-II observations.}
Retrieval-aware reliability distinguishes downstream-useful Bridges more clearly
than semantic similarity or local structural relevance, including within the Clean
candidate subset. This persistence after removing corrupted candidates suggests
that the model captures differences in downstream usefulness among plausible
Bridges, rather than merely separating clean from degraded inputs. Under matched
retrieval conditions, its advantage translates most consistently into stronger
mechanistic evidence recovery. Path success does not improve in parallel, however,
so better Bridge ranking should not be read as a uniform gain across retrieval
endpoints. The distinction is consistent with the supervision target: Stage II
learns which intermediate candidates contribute useful evidence, while complete
path recovery also depends on the fixed graph and search constraints. On EpiKG,
the utility signal remains informative despite a different biomedical schema and
a more imbalanced candidate distribution. Together, these results support
retrieval-aware reliability as a way to prioritize useful Bridges across settings,
with its clearest effect on the mechanistic evidence recovered by retrieval
under matched candidate access, search constraints, and guidance strength.

The single-Bridge comparison uses \rev{248 queries under five conditions (1,240 paired cases), not 1,240 independent queries.} \rev{Lexical $\rightarrow$ reliability PS/ER are Gene 31.87/36.46 $\rightarrow$ 31.38/38.82\%, Metabolite 25.71/33.47 $\rightarrow$ 25.71/36.78\%, and Process 28.46/41.17 $\rightarrow$ 27.31/41.87\%. ER increases in all tasks, while PS is unchanged or lower.} A positive $\Delta_h$ label can reward ER-only improvement; it does not guarantee endpoint completion. This $K=1$ intervention validates usefulness ranking, not the full $K=3$ selector. The complete-selector comparison below tests the larger-pool setting; implementation and automatic-candidate usefulness diagnostics are in Appendix~\ref{app:implementation}.

\begin{wraptable}{r}{0.47\textwidth}
\centering
\caption{Automatic-pool selection ($K=3$, $g=0.3$). Proposal is equal-weighted; the other rows share learned confidences.}
\label{tab:automatic_selector}
\small
\setlength{\tabcolsep}{3.5pt}
\begin{tabular}{@{}lrrr@{}}
\toprule
Selection & PS (\%) & ER (\%) & Utility \\
\midrule
Proposal Top3 & \rev{35.08} & \rev{38.73} & \rev{0.3618} \\
Confidence-only & \rev{32.66} & \rev{40.87} & \rev{0.3512} \\
Complete greedy & \rev{33.06} & \rev{41.77} & \rev{0.3568} \\
\bottomrule
\end{tabular}
\end{wraptable}

\noindent\textbf{Selection when the pool exceeds the budget.}
We test the original empty-start greedy rule on normalized automatic pools with $K=3$, $g=0.3$ and fixed search budgets. Appendix~\ref{app:implementation} specifies the Red/type terms and tie rules. TRAIN/VAL selection fixes $C=0.1$, $\mu=0.5$, $\beta=0.1$. \rev{Of 248 test pools, 239 exceed $K$, and greedy changes 101 selected sets relative to confidence-only.} The shared confidence scores isolate subset choice from confidence weighting.



\rev{Greedy minus confidence-only utility is $+0.00554$ (95\% CI $[-0.01420,0.02451]$). The complete automatic pipeline reaches PS/ER 33.47/41.84\%; its utility difference from Proposal Top3 is $-0.00197$ ($[-0.04268,0.03855]$). A validation-selected constant matches its per-query outcomes.} Appendix~\ref{app:implementation} details this combination. \rev{These reused-test comparisons bound whole-selector and gate claims; they do not negate the separate usefulness and type-complementarity evidence.}
% ============================================================

% 5.5 ADAPTIVE CONTROL OF SEMANTIC GUIDANCE
% ============================================================

\subsection{Adaptive Control of Semantic Guidance}
\label{sec:stage3}
\textbf{Query-specific influence control.} To answer RQ4, we now use the exact Table~\ref{tab:overall_gain} Clean Full configuration: identical candidates, confidences, features, predictions and search, with $g=0.5\sqrt{p_{\mathrm{rel}}}$. \rev{Its mean test gate is 0.36035. The validation-mean constant is 0.35877, and the validation-selected global constant is 0.45.} Five hundred seeded global permutations preserve the test gate distribution. The revised table corrects the previously reported \rev{ER 38.05\%/mean gate 0.10917} row with controls using the actual main Full gate vector; the prior summary is retained in the audit correction record.
\vspace{-0.15em}

\noindent\textbf{Control calibration and attribution.}
The controller uses the original 42-combination validation grid across five candidate conditions, with a nearest-0.05 cache during selection and continuous gates during test replay. In contrast, the global constant is selected on Clean validation utility, and the mean constant uses validation gates only. These tuning scopes are stated explicitly because their calibration distributions differ. All attribution arms then fix the same selected nodes, confidence scores, graph and search budgets. The 500 global shuffles reuse the frozen test gate vector without fitting new parameters. Thus their comparison tests query--gate correspondence while preserving the distribution of guidance strength; it does not evaluate a newly optimized controller. Paired resampling retains all task and candidate-condition variants of each mechanism record together; no control selects its parameters from test outcomes.

\begin{table}[H]
\centering
\caption{Stage-III controls using the actual Table~\ref{tab:overall_gain} Full configuration.}
\label{tab:stage3_controls}
\vspace{-0.35em}
\small
\begin{tabular*}{\textwidth}{@{\extracolsep{\fill}}lrrrr@{}}
\toprule
Method & PS (\%) & ER (\%) & Utility & Mean $g$ \\
\midrule
Graph Only & 27.02 & 33.97 & 0.2910 & 0.00000 \\
Validation-selected constant & \rev{39.11} & \rev{59.74} & \rev{0.4530} & \rev{0.45000} \\
Validation-mean constant & \rev{39.11} & \rev{59.54} & \rev{0.4524} & \rev{0.35877} \\
Shuffled Full (500-run mean) & \rev{39.24} & \rev{59.06} & \rev{0.4518} & \rev{0.36035} \\
Full (same as Table~\ref{tab:overall_gain}) & \rev{39.52} & \rev{59.46} & \rev{0.4550} & \rev{0.36035} \\
\bottomrule
\end{tabular*}
\end{table}

\vspace{-0.45em}
% ============================================================
% OA + HALF-WIDTH WRAPPED FIGURE 4
% ============================================================

\begin{figure}[H]

    \centering
    \setlength{\abovecaptionskip}{0pt}
    \setlength{\belowcaptionskip}{0pt}

    \includegraphics[
        width=1.0\linewidth,
        trim=0 22 0 4,
        clip,
        keepaspectratio
    ]{figure5.pdf}

    \caption{
Corrected main-configuration shuffle distributions (500 permutations). Full is marked by a vertical line; \rev{matching effects are modest and endpoint-dependent.}
    }
    \label{fig:stage3_shuffle}

\end{figure}
\vspace{-0.4em}

\textbf{Stage-III observations.} \rev{Full does not uniformly exceed strong constants: its ER is slightly lower, and Full minus the validation-mean constant has utility difference $0.00258$ (paired 95\% CI $[-0.00072,0.00861]$). Across 248 queries, utility improves in one, declines in one and is unchanged in 246. One-sided shuffle tests give ER/utility $p=0.0679/0.0220$; Holm adjustment across four endpoints gives $0.2036/0.0878$.} Figure~\ref{fig:stage3_shuffle} shows the corrected distributions. Query matching can affect paths, but \rev{these controls do not establish a stable independent gate advantage over strong constants.} Appendix~\ref{app:implementation} provides the positive/negative cases, complete configuration and uncertainty.

\noindent\textbf{Query-level benefit and harm.}
In Tolbutamide/type-2-diabetes Gene, Full at 0.28795 versus constant 0.35877 redirects a potassium-channel continuation toward an insulin-secretion path: \rev{PS changes $0\rightarrow1$, with ER unchanged at $1/7$}. Conversely, in Afamelanotide/erythropoietic-protoporphyria Gene, Full at 0.23878 reaches the target but \rev{ER falls from 1.0 to 0.8}. These graph paths illustrate benefit and harm, not additional biological validation. The first case changes path success without increasing reference-edge recovery; the second preserves success while losing reference edges. Reporting both endpoints exposes tradeoffs that either metric alone would conceal.

\noindent\textbf{Sensitivity to outgoing-edge order.}
With models and predictions fixed, original order, sorted order and ten seeded permutations give Full minus Graph Only mean \rev{PS gains of 9.68--13.71 and ER gains of 22.91--25.49 percentage points}. Full minus Semantic+Fixed remains positive in mean: \rev{PS 0.40--0.81 and ER 0.054--0.901 points}. These are paired point-difference ranges, not uniform significance. The original order is retained; Appendix~\ref{app:implementation} gives counting rules and all outcomes. With a fixed expansion budget, adjacency order can alter the explored frontier. Each comparison uses the same order across variants to isolate guidance under matched search.

\vspace{0.45em}
% ============================================================
% 6 CONCLUSION
% ============================================================
\vspace{-8pt}
\section{Conclusion}
We study how intermediate semantic information can guide biomedical
knowledge-graph retrieval while preserving graph-valid search. Our results
support a failure-oriented view of semantic guidance: useful intermediate
candidates must first be available, their downstream retrieval value must be
distinguished from semantic plausibility, and \rev{query-specific influence requires
matched fixed-strength controls}. Confidence-Gated Mechanistic Bridge Retrieval addresses these
coverage, utility, and influence failures through multi-view construction,
retrieval-aware Bridge reliability, and query-specific control, respectively.
\rev{Controlled analyses distinguish these retrieval behaviors: automatic proposal broadens useful-candidate access, and retrieval-derived scores discriminate usefulness beyond semantic similarity. The supplied-Clean gains remain conditional on candidate access, while complete greedy selection and query-specific control have not shown stable independent advantages over strong controls. The evidence therefore qualifies the performance claims while retaining the original decomposition of candidate coverage, downstream utility, and query-specific influence.}


% ============================================================
% AI USE STATEMENT
% ============================================================
\newpage
\subsection*{AI use statement}

Generative AI tools assisted literature search, manuscript editing, code development/debugging, and the implementation and execution of computational analyses during revision. Numerical results were computed by executable programs from the research data rather than invented by a language model. The authors remain responsible for verifying references, methods, analyses and the final manuscript.

% ============================================================
% ETHICS STATEMENT
% ============================================================
\subsection*{Ethics statement}

This work adheres to the ICLR Code of Ethics. The study does not 
involve human subjects, clinical interventions, or the collection of 
personally identifiable information. Our experiments use existing 
research datasets and biomedical knowledge resources in accordance 
with their intended research purposes. The proposed method is intended 
for research on biomedical knowledge retrieval and should not be 
interpreted as providing clinical or medical advice. We recognize that 
errors in biomedical knowledge retrieval may propagate incomplete or 
incorrect information; accordingly, the system is evaluated as a 
research method rather than as a tool for clinical decision-making.

% ============================================================
% REPRODUCIBILITY STATEMENT
% ============================================================
\subsection*{Reproducibility statement}

We provide the information necessary to reproduce the experiments 
reported in this work. The main paper specifies the proposed method, 
datasets, evaluation protocols, and experimental settings, while 
additional implementation details and configurations are documented 
in the Appendix. Controlled variants and evaluation procedures are 
described alongside the corresponding experiments to enable direct 
verification of the reported comparisons. These details are intended 
to support independent reproduction and verification of our findings.

\bibliography{iclr2027_conference}
\bibliographystyle{iclr2027_conference}

\appendix


% ============================================================
% A RELATED WORK
% ============================================================
\section{Related Work}
\label{app:related_work}

Knowledge-graph reasoning has progressed from embedding-based relational inference
and multi-hop path search~\citep{bordes2013transe,yang2017neurallp,das2018minerva,saxena2020embedkgqa}
toward retrieval systems that combine explicit graph structure with language-model
reasoning. Think-on-Graph and Think-on-Graph 2.0 perform iterative relation and
entity selection during graph traversal~\citep{sun2024tog,ma2025tog2}, whereas
GNN-RAG and SubgraphRAG learn query-conditioned relevance over graph elements
or subgraphs~\citep{mavromatis2025gnnrag,li2025subgraphrag}. Other recent
approaches emphasize graph-valid reasoning, relation-aware retrieval, and
structured medical evidence~\citep{luo2024rog,luo2025gcr,um2026rporag,wu2025medgraphrag}.
Biomedical knowledge graphs make these retrieval decisions especially important
because relevant evidence spans heterogeneous entity and relation types. Hetionet
and PrimeKG integrate broad biomedical knowledge
~\citep{himmelstein2017hetionet,chandak2023primekg}, DrugMechDB represents
curated mechanisms through ordered intermediate relations
~\citep{gonzalez2023drugmechdb}, and BioThings Explorer and BTE-RAG support
evidence-grounded biomedical retrieval
~\citep{callaghan2023biothings,joy2026bterag,matsumoto2024kragen}. In contrast to methods that
primarily determine which graph elements or regions should be retrieved, our work
studies when intermediate semantic hypotheses should influence an otherwise
graph-constrained retrieval process. Adaptive retrieval provides a complementary perspective. FLARE uses model
confidence to determine when additional evidence should be retrieved
~\citep{jiang2023flare}, Self-RAG learns reflection signals for adaptive retrieval
and evidence use~\citep{asai2024selfrag}, and Adaptive-RAG selects retrieval
strategies according to estimated query complexity
~\citep{jeong2024adaptiverag}. These methods establish that retrieval behavior
should vary across inputs, but they primarily adapt whether or how retrieval is
invoked. Our setting instead keeps graph-valid retrieval active and adapts the
influence of already selected semantic guidance. The distinction is the retrieval problem and controlled decomposition, not a claim that logistic regression, greedy diversification or score interpolation is itself novel. Self-RAG learns reflection signals for passage retrieval and generation; BTE-RAG supplies federated biomedical context for generated answers. Here known endpoints and a fixed search budget define intermediate path/evidence recovery, so its PS/ER is not directly comparable with generated-answer accuracy.
\small
\section{Implementation Details}
\label{app:implementation}

\paragraph{Retrieval configuration.}
All controlled variants operate on the same underlying graph representation and
use matched search constraints within each benchmark. Retrieval expands only
graph-valid outgoing transitions, and semantic guidance affects path ranking
without introducing virtual edges or modifying graph connectivity. For the
endpoint-conditioned DrugMechDB experiments, search is initialized from the
source entity and terminates when the target is reached or the search budget is
exhausted. We use beam size 2, a maximum of 8 hops, a 14-node budget, and a
20-edge expansion budget. Repeated nodes within a path are excluded. Candidate
paths are ranked after each expansion step, and only the highest-scoring beam
states are retained. The same search configuration is used across matched
DrugMechDB comparisons so that differences between variants arise from semantic
guidance rather than changes in search capacity.
The exploratory corruption study in Section~\ref{sec:exploratory} uses 1,000
retrieval instances and five random seeds under these matched constraints.

\noindent\textbf{Expansion order.} Outgoing $(\mathrm{neighbor},\mathrm{relation})$ transitions retain graph insertion order and first-occurrence deduplication. The edge counter increments before cycle and node-budget checks. With predictions fixed, original order, sorted order and ten seeded permutations yield Full minus Graph Only \rev{mean PS gains of 9.68--13.71 and ER gains of 22.91--25.49 percentage points. Full minus Semantic+Fixed means remain positive across these 12 orders (PS 0.40--0.81; ER 0.054--0.901 points), without establishing significance in every order.} The original main order is retained; no best-test-order selection is performed.

\noindent\textbf{Clean candidate construction.} Clean test inputs contain \rev{692 reference-member candidates over 248 queries}, including the intermediate answer in every pool. They supply node identities, not a complete path or relation sequence; graph retrieval remains necessary. A global entity-type map traverses reference records in file order, retaining the last label for each canonical ID. For each query, the constructor traverses its mechanism in stored node order, excludes known endpoints and deduplicates IDs. It places the reference intermediate answer first, preferentially adds nodes of unrepresented types, and fills remaining slots in record order to obtain up to three base candidates. Finally, it appends the first eligible reference node with the answer's type if absent. \rev{Five TRAIN lists therefore contain four nodes; all VAL/TEST lists contain at most three. The rule exactly reproduces all 1,354 TRAIN, 239 VAL and 248 TEST lists, including order, in both the training-sample and prediction tables (3,682 lists).} This establishes output-equivalent reconstruction, not authentication of the historical writer. The rule was identified using saved lists, including TEST, and is not independent validation. Clean comparisons test confidence weighting and control under reference-informed access, not automatic proposal or nontrivial subset filtering.

\noindent\textbf{Automatic candidate construction.} The proposer uses $0.35J_{\mathrm{token}}+0.20A_s+0.20A_t+0.10\mathrm{AntiHub}+0.075R_s+0.075R_t$, where $A_s,A_t$ indicate source/candidate and candidate/target adjacency and $R_s,R_t$ summarize endpoint relation quality. It excludes endpoints and reads neither intermediate answers nor query-specific reference membership. The cap of 600 is inactive (maximum pool 483). Figure~\ref{fig:stage1_multiview} and automatic retrieval share \rev{14,029} manifest-endpoint query--candidate records; the earlier 14,041-record raw-endpoint plot remains in the correction archive. Coverage uses the full view union; fixed-budget proposal retains the top three by this rule.

\noindent\textbf{Complete-selector operational test.} Selection follows the stated gain formula and empty-start $K=3$ greedy procedure. Red is maximum one-hop undirected-neighbor Jaccard with an already selected node (zero for an empty set); TypeBonus indicates a new raw graph type, and node-ID order breaks ties. These definitions specify the selector evaluated here. TRAIN-only preprocessing and class-balanced LR use $C\in\{0.1,1,10\}$; a VAL grid $\mu\in\{0,0.1,0.25,0.5\}$, $\beta\in\{0,0.05,0.1\}$ selects $C=0.1,\mu=0.5,\beta=0.1$ before reused-test analysis. \rev{There are 239/248 oversized test pools, and greedy differs from confidence-only in 101 selections. Greedy PS/ER is 33.06/41.77\%, versus confidence-only 32.66/40.87\%; the utility difference $0.00554$ has paired interval $[-0.01420,0.02451]$. Against proposal Top3, PS is lower and ER higher. This does not establish stable gains from the full rule. On 14,029 automatic test candidates, the same model's singleton-usefulness AUROC/AUPRC is 0.870/0.108, versus lexical 0.609/0.036 and graph heuristic 0.346/0.012 (prevalence 1.65\%); discrimination is not proof of better final subset retrieval.} These diagnostics reuse an already inspected test set and do not provide fresh independent confirmation.

\noindent\textbf{Complete automatic pipeline.} Automatic proposal, the Table~\ref{tab:automatic_selector} greedy selector and query-specific control are evaluated together with the same pools, learned confidences, $C=0.1$, $\mu=0.5$, $\beta=0.1$, $K=3$ and search budgets. The existing TRAIN-only gate checkpoint and $g=0.5\sqrt{p_{\mathrm{rel}}}$ mapping remain frozen; original query features are recomputed on the selected Bridges. No model is refitted. On the same greedy-selected VAL inputs, utility selects $g=0.45$ from $\{0,0.05,\ldots,0.5\}$ before TEST evaluation, with ties at 12 decimal places resolved toward the smaller constant.

\begin{table}[t]
\centering
\caption{Complete automatic pipeline and matched guidance controls. All rows share the graph and search budgets; the last two share selected nodes and confidences.}
\label{tab:automatic_pipeline}
\small
\setlength{\tabcolsep}{5pt}
\begin{tabular}{lrrr}
\toprule
Automatic configuration & PS (\%) & ER (\%) & Utility \\
\midrule
{\color{blue}Proposal Top3, equal confidence, $g=0.3$} & \rev{35.08} & \rev{38.73} & \rev{0.36176} \\
{\color{blue}Confidence-only, learned confidence, $g=0.3$} & \rev{32.66} & \rev{40.87} & \rev{0.35123} \\
{\color{blue}Complete greedy, learned confidence, $g=0.3$} & \rev{33.06} & \rev{41.77} & \rev{0.35677} \\
{\color{blue}Complete greedy, query-specific gate} & \rev{33.47} & \rev{41.84} & \rev{0.35980} \\
{\color{blue}Same greedy, validation-selected $g=0.45$} & \rev{33.47} & \rev{41.84} & \rev{0.35980} \\
\bottomrule
\end{tabular}
\end{table}

Table~\ref{tab:automatic_pipeline} evaluates the complete combination directly. \rev{The gate's mean TEST strength is 0.408227. Its PS, ER and utility equal the selected constant for every query; the paired utility difference and interval are both zero. Relative to Proposal Top3, utility differs by $-0.001965$ (95\% CI $[-0.042682,0.038555]$), with PS $-1.61$ points ($[-6.32,3.11]$) and ER $+3.11$ points ($[-0.45,6.61]$).} Intervals use 5,000 mechanism-grouped paired bootstrap draws (seed 20260914), without multiplicity adjustment. This previously inspected TEST set is not independent confirmation. \rev{The comparison establishes neither overall superiority nor inferiority to Proposal Top3, and supplies no incremental gate advantage over the matched constant. These bounds do not negate the separately measured usefulness and type-complementarity results.}

\noindent\textbf{Selector term ablations.} Under the same frozen scores, $K=3$, $g=0.3$ and search budget, removing only TypeBonus gives \rev{PS/ER/utility 31.05/41.19\%/0.3409}; removing only Red gives \rev{33.06/41.80\%/0.3569}, compared with complete greedy \rev{33.06/41.77\%/0.3568}. Complete minus no-Type yields \rev{$+2.02$ PS points (paired pointwise 95\% CI $[0.40,4.02]$) and utility $+0.01588$ ($[0.00375,0.03138]$)}. Complete minus no-Red utility is \rev{$-0.00009$ ($[-0.01349,0.01301]$)}. Thus \rev{type complementarity has support in this component ablation, while an independent Red benefit and complete-rule superiority to confidence-only remain unestablished}. These are pointwise intervals from existing test records, not multiplicity-adjusted evidence for every term; no term is reselected using test outcomes.

In the existing automatic TRAIN singleton labels, \rev{973 of 1,313 positive cases improve ER without increasing PS, while 340 increase PS}. Both receive the same positive binary target. This documents why a usefulness label can reward partial evidence without path completion; it is a label diagnostic, not proof of the unique cause of final-set errors.

\noindent\textbf{Matched learned-path baseline.} One numeric logistic path head uses the same graph, Bridge confidences and search budgets, with \rev{77 features} from existing graph/Bridge score components and query interactions. There are \rev{101,124 admitted prefixes from 6,770 TRAIN query-condition cases}. Its labels compare prefix utility with Graph Only using the same reference source and $M=0.7\mathrm{PS}+0.3\mathrm{ER}$, but at path rather than query granularity. Query balancing and feature standardization use TRAIN only. The head fixes the original gate configuration ($C=0.5$, class-balanced LR, liblinear, seed 20260827). \rev{Clean Full minus this baseline has utility $+0.03582$ (95\% CI $[0.01602,0.05860]$). On fixed automatic proposal Top3, Full/path-scorer PS/ER are 34.68/39.12\% and 33.06/39.82\%, with utility interval spanning zero.} This single compact control does not establish superiority over neural or text-rich path scorers.

\noindent\textbf{Canonical gate controls and cases.} The corrected Table~\ref{tab:stage3_controls} uses exactly the main Clean fixtures, learned confidences and $g=0.5\sqrt{p_{\mathrm{rel}}}$. The original 42-combination controller validation uses all five candidate conditions, with a nearest-0.05 gate cache for selection and continuous gates for test replay; constants use Clean validation utility, and the mean constant uses validation gates only. \rev{Full minus validation-mean ER is $-0.0806$ points (95\% CI $[-0.2490,0.0000]$); utility is $+0.00258$ ($[-0.00072,0.00861]$).} Five hundred fixed global shuffles retain the test gate vector; the paired interval against per-query shuffle averages and the batch randomization test are distinct analyses. \rev{All 124,992 control outcomes have been replayed.} The archived \rev{ER 38.05\%/mean-gate 0.10917} result remains an explicitly corrected historical summary, not evidence for this configuration. In Tolbutamide/type-2-diabetes Gene, Full's lower gate 0.28795 versus constant 0.35877 changes a potassium-channel continuation to an insulin-secretion path reaching the target (\rev{PS $0\rightarrow1$, ER unchanged at $1/7$}). In Afamelanotide/erythropoietic-protoporphyria Gene, Full at 0.23878 reaches the target but \rev{ER falls from 1.0 to 0.8}. These illustrate observed graph paths, not additional biological-validity claims.

\noindent\textbf{Split and inference scope.} The original 1,841 queries have 1,540 mechanism IDs, with no ID crossing TRAIN/VAL/TEST. All candidate/corruption variants follow the same ID split. \rev{Fifteen question-text and fifteen drug--disease-pair groups span splits through different IDs; eight test queries have pairs in TRAIN, nine in TRAIN or VAL.} Excluding those nine as a post-hoc sensitivity leaves 239 queries: \rev{Clean Full minus the matched path scorer retains positive PS/ER/utility intervals $[0.85,6.15]$/$[0.26,4.02]$ points/$[0.01171,0.05126]$. This is a sensitivity subset, not an independent holdout.} Reference information is used for training labels and offline evaluation; it is not an automatic proposer input. The fixed aggregate graph remains transductive.

\paragraph{Scorer values.}
The fixed settings are $\lambda=0.82$, $C(p,H^*)=0.45C_{\mathrm{cov}}+0.25C_{\mathrm{mean}}+0.15C_{\mathrm{cur}}+0.15C_{\mathrm{recent}}$, and $\rho(p)\in\{0.65,0.85,1.0\}$ in ascending order of distinct Bridge hits.


\paragraph{Reliability estimators and feature preprocessing.}
The Bridge-level reliability estimator in Stage II uses the semantic and structural
features defined in Section~\ref{sec:bridge_reliability}. Semantic features are
constructed with word-level TF--IDF using 1--2 grams and character-level
TF--IDF using 3--5 grams. Both vocabularies are fitted on the training split only.
Numeric graph features are standardized using statistics estimated from the
training split. Bridge reliability is learned with a class-balanced logistic
regression model. The query-level estimator in Stage III follows the same
train-only preprocessing principle. Its textual representation combines the task,
query, and Bridge-set context, while its numeric representation contains only
pre-search graph and candidate-set statistics. Neither estimator uses
reference-path membership, shortest-path oracle information, or held-out
post-retrieval outcomes as input features.

\paragraph{Model selection and continuous-gate tuning.}
All learned models and controller parameters are selected using training and
validation data only. For the continuous gate, we evaluate
$g_{\max}\in\{0.20,0.25,0.30,0.35,0.40,0.45,0.50\}$ and
$\gamma\in\{0.5,0.75,1.0,1.25,1.5,2.0\}$ on the validation split. Candidate
settings are compared using the same retrieval utility used by the method,
$0.7\,\mathrm{PathSuccess}+0.3\,\mathrm{EdgeRecall}$, and the selected pair is
then fixed for held-out evaluation. Revision diagnostics reuse previously inspected held-out records and are explicitly labeled as such; they do not constitute a new blind test. No failed local optimization is substituted into the main method or tables. The constrained search is discrete and is
not differentiated through, and no downstream answer-generation loss is used
to train either reliability estimator or the retrieval controller. Held-out test
outcomes are never used for model selection, feature fitting, or hyperparameter
tuning.

\section{Sensitivity of Utility and Search Coefficients}
\label{app:search_parameter_sensitivity}

\paragraph{Stage-II utility weight.}
We recompute usefulness labels from recorded candidate-level retrieval outcomes for $\omega\in\{0.5,0.6,0.7,0.8,0.9\}$ without refitting the reliability model. Across 3,460 held-out candidates, including 692 Clean candidates, positive rates and frozen-score AUROC/AUPRC are unchanged: 0.795/0.393 overall and 0.815/0.655 on Clean. This post hoc check establishes local label stability, not a uniquely optimal weight or a full retrieval rerun.

\paragraph{Protocol.}
We replay DrugMechDB Clean with the main experiment's graph, Bridge confidences, query gates, and search budgets. The default replay reproduces all recorded outcomes for the 248 test queries. We vary one score component at a time without retraining or test-set selection. Table~\ref{tab:search_parameter_stress} reports Path Success (PS), Edge Recall (ER), and paths changed from the default on 239 validation and 248 test queries.

\begin{table}[H]
\centering
\small
\caption{Scoring-component interventions on DrugMechDB Clean. PS and ER are percentages; changed paths are counts. Interventions are post hoc, not hyperparameter selection.}
\label{tab:search_parameter_stress}
\begin{tabular}{lccc}
\toprule
Configuration & Validation PS / ER & Test PS / ER & Changed paths (val / test) \\
\midrule
Default & 29.71 / 54.28 & 39.52 / 59.46 & 0 / 0 \\
No relation prior ($\lambda=0$) & 31.38 / 50.63 & 38.31 / 55.20 & 107 / 97 \\
No hub term ($\lambda=1$) & 28.45 / 53.41 & 37.10 / 59.17 & 49 / 38 \\
No one-Bridge reward ($\rho_1=0$) & 29.71 / 50.40 & 37.10 / 54.22 & 24 / 29 \\
No two-Bridge reward ($\rho_2=0$) & 23.01 / 35.41 & 31.05 / 39.11 & 138 / 152 \\
Mean-confidence term only & 28.03 / 48.48 & 35.48 / 51.85 & 49 / 66 \\
\bottomrule
\end{tabular}
\end{table}

\paragraph{Interpretation.}
Removing relation quality or Bridge rewards changes retrieval, most strongly when the two-Bridge reward is removed. Nearby coefficients are less identifiable: changing $0.82$ to $0.819$, $0.81$, or $0.83$ changes no path on either split. One-at-a-time $\pm0.01$ changes to Bridge weights or the first two saturation factors leave PS and ER unchanged; $\pm0.05$ changes shift ER by at most 0.084 validation and 0.081 test percentage points. Thus the components matter, but no default coefficient is uniquely optimal.

The original $0.012|p|$ penalty is rank-inert: paths within each beam step have equal length, and search stops upon reaching the target. Replays at zero and $0.2$ change no path; Section~\ref{sec:confidence_gate} states the effective score. This frozen-search analysis excludes retraining and graph transfer.
\normalsize

\section{Transfer Beyond DrugMechDB}
\label{app:transfer}

\paragraph{Scope and protocol.}
We examine whether graph-valid retrieval with semantic Bridge guidance remains useful when the graph or answer format changes. These experiments test transfer of a retrieval principle, rather than transfer of DrugMechDB's biomedical mechanism labels or its path-success metric. BioHopR~\citep{kim2025biohopr} uses a PrimeKG-based, target-agnostic two-hop entity-ranking task; ComplexWebQuestions (CWQ)~\citep{talmor2018web} uses question-specific Freebase subgraphs and target-agnostic multi-hop answer ranking. Gold answers and paths are excluded from test-time scoring in both settings. Within each benchmark, variants share the same graph and candidate access, and model selection uses training and validation data. Hits@K and MRR are reported as percentages, so only within-benchmark comparisons are meaningful.

BioHopR has 1,146 held-out test questions (seed 17); CWQ has 3,531, and its table entries average seeds 17, 29, and 43.

\begin{table}[t]
\centering
\small
\setlength{\tabcolsep}{4pt}
\caption{Transfer retrieval under matched within-benchmark comparisons. Scores are percentages; gains over Graph Only are percentage points. Bold marks the best score within each benchmark.}
\label{tab:transfer_results}
\begin{tabular*}{\linewidth}{@{\extracolsep{\fill}}llrrrrr@{}}
\toprule
Benchmark & Retrieval variant & Hits@1 & Hits@5 & MRR & $\Delta$H@1 & $\Delta$MRR \\
\midrule
BioHopR & Graph Only & 81.33 & 89.88 & 85.30 & --- & --- \\
        & Semantic Only & \textbf{82.29} & 89.18 & 85.84 & +0.96 & +0.54 \\
        & Full & \textbf{82.29} & \textbf{90.14} & \textbf{86.06} & +0.96 & +0.76 \\
\addlinespace
CWQ & Graph Only & 0.03 & 0.14 & 0.51 & --- & --- \\
    & Semantic Only & 0.25 & 2.10 & 1.68 & +0.23 & +1.17 \\
    & Fixed Bridge & \textbf{4.64} & \textbf{12.02} & \textbf{8.24} & +4.62 & +7.73 \\
    & Full target-agnostic & 4.47 & 10.95 & 7.75 & +4.44 & +7.24 \\
\bottomrule
\end{tabular*}
\end{table}

\paragraph{What transfers.}
On BioHopR, Full improves Hits@1 by 0.96 points and MRR by 0.76 points over structural retrieval. This is evidence that semantic guidance can contribute on a different biomedical graph and an answer-ranking task, but the gain is modest and does not establish a broad mechanistic-reasoning benefit. On CWQ, Bridge-guided retrieval substantially improves over Graph Only despite the shift to general-domain Freebase reasoning: mean Hits@1 rises from 0.03\% to 4.47\%. A paired bootstrap on seed 17 gives a 95\% interval of [3.82, 5.21] percentage points for Full minus Graph Only. However, the fixed-Bridge variant is stronger than Full on all three reported CWQ metrics; the corresponding Hits@1 interval for Full minus Fixed Bridge is [-0.31, 0.08] points. Thus CWQ supports transfer of graph-constrained Bridge guidance, not a consistent advantage for learned query-specific influence. The low absolute CWQ accuracy further limits any claim of competitive general-domain question answering.

\paragraph{OpenBioLink and limits of the evidence.}
OpenBioLink~\citep{breit2020openbiolink} changes the biomedical graph schema and relation distribution and would provide a complementary graph-transfer test. We do not report a numerical OpenBioLink comparison here because a matched run and its per-query outputs are not available for verification. Consequently, the present appendix does not establish robustness across all three proposed transfer settings. The BioHopR and CWQ results support a narrower conclusion: semantic Bridges can remain useful after a change in graph and task format, whereas the incremental value of reliability-controlled influence depends on the target setting. Neither establishes mechanistic-path recovery across arbitrary graphs or answer-ranking tasks.
\normalsize

\begingroup
\makeatletter
\setlength{\@fptop}{0pt}
\makeatother
\clearpage
\endgroup
\end{document}
